\chapter{First Steps}
\index{expressions}
\label{ch:first-steps}

With Kaappi installed and the REPL ready, it is time to write your first Scheme code. This chapter covers the fundamental building blocks: expressions, basic arithmetic, variables, and defining simple functions. Open a REPL and type along---every example here is meant to be tried interactively.

\section{Everything Is an Expression}

In Python, you distinguish between \emph{statements} (which do things) and \emph{expressions} (which produce values). In Scheme, there is only one category: everything is an expression, and every expression produces a value.

The syntax is uniform. An expression is either an atom (a number, string, or symbol standing alone) or a list enclosed in parentheses\index{S-expression}. When the evaluator sees a list, it treats the first element as the operator and the rest as operands:

\begin{lstlisting}[style=repl]
> (+ 2 3)
5
> (* 6 7)
42
> (- 100 1)
99
\end{lstlisting}

The operator always comes first. This is called \emph{prefix notation}\index{prefix notation}. There is no infix \texttt{2 + 3} in Scheme---it is always \texttt{(+ 2 3)}.

\subsection{No Precedence Rules}

In Python, \texttt{2 + 3 * 4} evaluates to \texttt{14} because \texttt{*} binds tighter than \texttt{+}. In Scheme, you nest expressions explicitly:

\begin{lstlisting}[style=repl]
> (+ 2 (* 3 4))
14
> (* (+ 1 2) (- 10 4))
18
\end{lstlisting}

The parentheses \emph{are} the precedence. There is nothing to memorize, and there is never ambiguity about evaluation order.

\begin{note}[Coming from Python]
If \texttt{(+ 2 3)} looks strange, think of it as a procedure\index{procedure} call where the procedure name goes inside the parentheses: like \texttt{add(2, 3)} but without the commas and with the procedure name moved inside. After a few examples, the pattern becomes natural.
\end{note}

\subsection{Operators Take Multiple Arguments}

Unlike most languages where \texttt{+} is binary, Scheme's arithmetic operators accept any number of arguments:

\begin{lstlisting}[style=repl]
> (+ 1 2 3 4 5)
15
> (* 2 3 4)
24
> (+)
0
> (*)
1
\end{lstlisting}

With zero arguments, \texttt{+} returns 0 (the additive identity) and \texttt{*} returns 1 (the multiplicative identity).

\section{Atoms}
\index{atoms}

Atoms are the simplest values in Scheme---they cannot be broken into smaller parts.

\subsection{Numbers}
\index{numbers}

Kaappi supports a full numeric tower:

\begin{lstlisting}[style=repl]
> 42
42
> -7
-7
> 3.14
3.14
> 1/3
1/3
> (+ 1/3 1/6)
1/2
\end{lstlisting}

Integers and exact rationals stay exact---there is no floating-point drift when you work with fractions. Kaappi also supports arbitrarily large integers (bignums):

\begin{lstlisting}[style=repl]
> (* 1000000000 1000000000 1000000000)
1000000000000000000000000000
\end{lstlisting}

\begin{note}[Coming from Python]
Python 3 also has arbitrary-precision integers, so this will feel familiar. The difference is that Scheme has exact rationals (\texttt{1/3}) as a built-in type---Python requires the \texttt{fractions} module.
\end{note}

\subsection{Strings}
\index{strings}

Strings are delimited by double quotes:\index{string-length}\index{string-append}

\begin{lstlisting}[style=repl]
> "hello, world"
"hello, world"
> (string-length "hello")
5
> (string-append "foo" "bar")
"foobar"
\end{lstlisting}

Single quotes mean something entirely different in Scheme (they denote quotation, which we will cover shortly), so strings are always double-quoted.

\subsection{Booleans}
\index{booleans}

The boolean values are \texttt{\#t} (true) and \texttt{\#f} (false). There is one critical rule that surprises Python programmers: \textbf{only \texttt{\#f} is false}. Everything else---including \texttt{0}, the empty string, and the empty list---is true:

\begin{lstlisting}[style=repl]
> (if 0 "truthy" "falsy")
"truthy"
> (if "" "truthy" "falsy")
"truthy"
> (if '() "truthy" "falsy")
"truthy"
> (if #f "truthy" "falsy")
"falsy"
\end{lstlisting}

\begin{warning}[Only \texttt{\#f} is false]
This is one of the most common sources of bugs when coming from Python. In Scheme, \texttt{(if 0 ...)} takes the ``then'' branch because \texttt{0} is not \texttt{\#f}. If you want to test whether a number is zero, use \texttt{(zero? n)} or \texttt{(= n 0)} explicitly.
\end{warning}

\subsection{Characters}
\index{characters}

Individual characters are written with a \texttt{\#\textbackslash} prefix:

\begin{lstlisting}[style=repl]
> #\a
#\a
> #\space
#\space
> #\newline
#\newline
> (char-alphabetic? #\Z)
#t
\end{lstlisting}

Characters are distinct from single-character strings.

\subsection{Symbols}
\index{symbols}

A symbol is a lightweight, interned name. Symbols are compared by identity (fast pointer comparison), not by character content like strings:

\begin{lstlisting}[style=repl]
> 'hello
hello
> (symbol? 'hello)
#t
> (eq? 'abc 'abc)
#t
\end{lstlisting}

The \texttt{'} (single quote) is called \emph{quote}\index{quote}. It prevents the expression that follows from being evaluated. Without it, \texttt{hello} would be looked up as a variable name, causing an ``unbound variable'' error.

\section{Defining Variables}
\index{define}

Use \texttt{define} to bind a name to a value:

\begin{lstlisting}[style=repl]
> (define pi 3.14159)
> pi
3.14159
> (* pi 10 10)
314.159
\end{lstlisting}

Once defined, you can use the name anywhere an expression is expected. Names in Scheme can contain characters that would be illegal in most languages: hyphens, question marks, and exclamation marks are all common:

\begin{lstlisting}[style=repl]
> (define my-favorite-number 42)
> my-favorite-number
42
\end{lstlisting}

\begin{note}[Naming conventions]
Scheme uses \texttt{kebab-case}\index{kebab-case} (words separated by hyphens) instead of \texttt{snake\_case} or \texttt{camelCase}. Predicates end with \texttt{?} (e.g., \texttt{zero?}, \texttt{string?}). Mutating procedures end with \texttt{!} (e.g., \texttt{set!}, \texttt{vector-set!}).
\end{note}

\section{Defining Functions}
\index{functions}\index{lambda}\index{procedures}

Scheme officially calls these \emph{procedures}\index{procedure} rather than functions, because they can have side effects (like printing output). This book uses both terms interchangeably.

To define a procedure, put the name and parameters in a list after \texttt{define}:

\begin{lstlisting}[style=repl]
> (define (square x) (* x x))
> (square 7)
49
> (square 1.5)
2.25
\end{lstlisting}

This is shorthand for the more explicit form using \texttt{lambda}, which creates an anonymous function:

\begin{lstlisting}[style=repl]
> (define square (lambda (x) (* x x)))
> (square 5)
25
\end{lstlisting}

\texttt{lambda}\index{lambda} is like Python's \texttt{lambda} keyword, but without the one-expression limitation---a Scheme lambda can contain any number of expressions in its body.

Functions can take multiple parameters:

\begin{lstlisting}[style=repl]
> (define (area width height) (* width height))
> (area 5 3)
15
\end{lstlisting}

And they can call other functions, including themselves:

\begin{lstlisting}[style=repl]
> (define (sum-of-squares a b)
    (+ (square a) (square b)))
> (sum-of-squares 3 4)
25
\end{lstlisting}

\subsection{Functions Are Values}

Functions are ordinary values in Scheme. You can pass them as arguments to other functions:

\begin{lstlisting}[style=repl]
> (define (apply-twice f x) (f (f x)))
> (apply-twice square 3)
81
> (apply-twice square 2)
16
\end{lstlisting}

In the first call, \texttt{(square 3)} gives \texttt{9}, then \texttt{(square 9)} gives \texttt{81}. In the second, \texttt{(square 2)} gives \texttt{4}, then \texttt{(square 4)} gives \texttt{16}. This ability to pass functions around is what makes Scheme a \emph{functional} language, and it is the foundation for many powerful patterns you will learn in later chapters.

\section{Basic Conditionals}
\index{if}\index{conditionals}

The \texttt{if} expression has three parts: a test, a ``then'' expression, and an ``else'' expression:

\begin{lstlisting}[style=repl]
> (if (> 5 3) "yes" "no")
"yes"
> (if (< 5 3) "yes" "no")
"no"
\end{lstlisting}

Because \texttt{if} is an expression (not a statement), it produces a value. You can use it anywhere a value is expected:

\begin{lstlisting}[style=repl]
> (define (abs x)
    (if (< x 0) (- x) x))
> (abs -7)
7
> (abs 3)
3
\end{lstlisting}

\begin{note}[Coming from Python]
This is like Python's ternary expression \texttt{"yes" if 5 > 3 else "no"}, but in Scheme the condition comes first. The pattern is always \texttt{(if test then else)}.
\end{note}

For multiple branches, use \texttt{cond}\index{cond}. It works like Python's \texttt{if/elif/else}:

\begin{lstlisting}[style=repl]
> (define (classify n)
    (cond
      ((< n 0) "negative")
      ((= n 0) "zero")
      (else "positive")))
> (classify -5)
"negative"
> (classify 0)
"zero"
> (classify 42)
"positive"
\end{lstlisting}

Each clause is \texttt{(test result)}. The \texttt{else} clause matches when none of the earlier tests are true.

\section{Comments}
\index{comments}

Line comments start with a semicolon and extend to the end of the line:

\begin{lstlisting}
; This is a comment
(define x 10) ; inline comment
\end{lstlisting}

Block comments use \texttt{\#|} and \texttt{|\#}:

\begin{lstlisting}
#|
This is a
multi-line comment
|#
\end{lstlisting}

A datum comment\index{datum comment} \texttt{\#;} comments out the next expression:

\begin{lstlisting}[style=repl]
> (+ 1 #;(this is ignored) 2)
3
\end{lstlisting}

This last form is particularly useful for temporarily disabling an argument or clause without deleting it.

\section{Putting It Together}

Let's write a slightly larger example that uses what we have learned. Here is a function that converts a temperature from Fahrenheit to Celsius:

\begin{lstlisting}[style=repl]
> (define (fahrenheit->celsius f)
    (* (/ 5 9) (- f 32)))
> (fahrenheit->celsius 212)
100
> (fahrenheit->celsius 32)
0
> (fahrenheit->celsius 72)
200/9
\end{lstlisting}

Notice that \texttt{(fahrenheit->celsius 72)} returns the exact rational \texttt{200/9} rather than a floating-point approximation. Kaappi preserves exactness when all inputs are exact. If you want a decimal, use a floating-point literal:

\begin{lstlisting}[style=repl]
> (define (fahrenheit->celsius-approx f)
    (* (/ 5.0 9.0) (- f 32)))
> (fahrenheit->celsius-approx 72)
22.22222222222222
\end{lstlisting}

\section*{Summary}

\begin{itemize}
    \item Scheme programs are built entirely from expressions in prefix notation, so there are no operator-precedence rules to memorize.
    \item The atomic data types are numbers, strings, booleans, characters, and symbols.
    \item \texttt{define} binds both variables and named functions, and functions are themselves first-class values.
    \item \texttt{if} and \texttt{cond} select between alternatives based on a condition.
    \item Code can be documented with line comments (\texttt{;}) and block comments.
\end{itemize}

\section*{Exercises}

\begin{exercise}[Predict and Check]
Before typing each expression into the REPL, predict what the result will be. Then check your answer:

\begin{lstlisting}[style=repl]
> (+ (* 3 5) (- 10 6))
> (/ 1 3)
> (+ 1/4 1/4 1/4 1/4)
> (* (+ 1 2) (+ 3 4) (+ 5 6))
> (- (* 2 3) (* 4 5))
\end{lstlisting}

Which results surprised you? Pay attention to how Kaappi handles exact rationals versus integers.
\end{exercise}

\begin{exercise}[Celsius to Fahrenheit]
The chapter showed a \code{fahrenheit->celsius} function. Write the inverse: a function \code{celsius->fahrenheit} that takes a temperature in Celsius and returns Fahrenheit. The formula is $F = C \times \frac{9}{5} + 32$. Test it in the REPL:

\begin{lstlisting}[style=repl]
> (celsius->fahrenheit 0)
32
> (celsius->fahrenheit 100)
212
> (celsius->fahrenheit 37)
493/5
\end{lstlisting}

In Python, this would be \code{def celsius\_to\_fahrenheit(c): return c * 9/5 + 32}.
\end{exercise}

\begin{exercise}[Absolute Value]
Write a function \code{(absolute-value n)} that returns the absolute value of \code{n} using \code{if}. It should handle positive numbers, negative numbers, and zero:

\begin{lstlisting}[style=repl]
> (absolute-value -7)
7
> (absolute-value 5)
5
> (absolute-value 0)
0
\end{lstlisting}

Compare your solution to the \code{abs} function defined earlier in this chapter. Are they equivalent?
\end{exercise}

\begin{exercise}[Recursive Power]
Write a function \code{(power base exponent)} that computes $\mathit{base}^{\mathit{exponent}}$ using recursion. Assume \code{exponent} is a non-negative integer. The base case is: anything raised to the power 0 is 1. The recursive case is: $b^n = b \times b^{n-1}$.

\begin{lstlisting}[style=repl]
> (power 2 10)
1024
> (power 3 0)
1
> (power 5 3)
125
\end{lstlisting}

In Python, you would write:

\begin{lstlisting}[language=Python, style=scheme, keywordstyle=\bfseries]
def power(base, exponent):
    if exponent == 0:
        return 1
    return base * power(base, exponent - 1)
\end{lstlisting}

Translate this logic into Scheme using \code{if}, \code{=}, and \code{-}.
\end{exercise}

\begin{exercise}[Temperature Classifier]
Write a function \code{(describe-temp t)} that takes a temperature in Celsius and returns a string describing it. Use \code{cond} with these ranges:

\begin{itemize}
    \item Below 0: \code{"freezing"}
    \item 0 to 15: \code{"cold"}
    \item 16 to 25: \code{"warm"}
    \item Above 25: \code{"hot"}
\end{itemize}

\begin{lstlisting}[style=repl]
> (describe-temp -5)
"freezing"
> (describe-temp 10)
"cold"
> (describe-temp 22)
"warm"
> (describe-temp 35)
"hot"
\end{lstlisting}

Hint: you will need both \code{<} and \code{<=} (or \code{>=}) to express the boundary conditions.
\end{exercise}

In the next chapter, we will explore Scheme's data types in depth.

\chapterend
